%% Conclusion

This study demonstrates that the economic feasibility and decarbonization of forest-residue-derived SAF depend on jointly optimizing plant scale, biomass availability, hydrogen supply and carbon capture strategies. Scale-up from 0.04 to 1.24 MMt/yr SAF production rate reduces MJSP by approximately 60\%, bringing it below \$1/L. However, approximately 90\% of the total cost reduction is achieved at 0.41 MMt/yr production scale, where economies of scale and the increasing cost of biomass supply intersect. Beyond this scale, single-train capacity limitations and feedstock price escalation offset further economies of scale, highlighting the need to balance market substitution potential against domestic forest-residue availability.

The grey-hydrogen base case exhibits GHG intensity of 100.8 g CO$_2$e/MJ, exceeding the CORSIA threshold of 75.93 g CO$_2$e/MJ. Integrating green ammonia-derived hydrogen and post-combustion CCS reduces emissions to as low as 21.9 g CO$_2$e/MJ, although the MJSP increases from \$0.79/L to \$3.09/L across the evaluated design space. CCS provides the more cost-effective decarbonization pathway per unit of GHG emissions avoided, making the 60\% CCS configuration the lowest-cost option for achieving CORSIA compliance. 

Green ammonia-derived hydrogen becomes economically competitive with CCS when the grey-hydrogen price exceeds a crossover point of \$1.59-2.91/kg H$_2$. This behavior reflects the strong influence of hydrogen supply costs on the MJSP. The economic advantage of green ammonia therefore depends primarily on its cost relative to grey hydrogen. The crossover price provides a simple benchmark for evaluating when green ammonia becomes more competitive than CCS.

This relationship provides a transparent basis for evaluating how future changes in hydrogen prices may alter the preferred decarbonization strategy. Overall, the findings establish a scale-dependent and cost-constrained framework for SAF deployment that integrates finite biomass resources with alternative decarbonization pathways. The results underscore that neither maximum plant scale nor maximum emission reduction alone defines the preferred configuration; instead, deployment decisions must jointly account for feedstock availability, production economics, and emissions compliance.

